\chapter{相关技术与理论基础}
\label{cha:related-technology}

\section{车联网与入侵检测概述}
\label{sec:related-vanet-ids}

车联网是由车载终端、路侧单元、边缘节点、云平台以及多类通信协议共同组成的开放式协同网络，其核心特征在于高移动性、强时变性和多主体异构协作。与传统企业网络相比，车联网中的节点地理分布更广，通信链路质量波动更明显，设备侧资源约束也更强。一旦攻击者利用伪造身份、恶意报文、重放数据或异常控制指令破坏通信流程，不仅会影响信息服务质量，还可能干扰车辆控制决策，因此安全防护必须兼顾实时性、鲁棒性和低开销\cite{placeholder_iov_security_survey,placeholder_vanet_ids_survey}。

入侵检测系统的任务是在海量运行数据中识别正常行为与异常行为，并尽可能降低漏报与误报。现有车联网入侵检测方案通常包括基于规则、基于统计分析以及基于机器学习三类思路。规则法适合识别已知模式，统计法关注行为偏离程度，而机器学习方法更擅长从复杂特征中学习潜在模式。对于当前项目而言，检测对象主要体现为表格化的车联网日志或网络流量记录，系统通过对原始字段执行缺失值处理、标签规范化与特征提取，再利用深度模型完成分类判别。这种“数据驱动检测”范式更适合在复杂环境下扩展，但其前提是能够在保证隐私和通信代价可控的前提下完成模型训练。

在车联网环境中，入侵检测还呈现出几个典型难点。第一，数据分布具有明显异构性，不同车辆、不同道路环境下采集到的样本差异较大；第二，车辆节点参与训练的时间窗口有限，且算力、电池状态与信道质量并不稳定；第三，原始数据中可能包含敏感位置信息、控制状态或行为轨迹，不适合直接集中上传。由此可见，车联网入侵检测不仅是模型分类问题，更是一个涉及分布式计算、资源调度与工程部署约束的系统问题。

\section{联邦学习基本原理}
\label{sec:related-federated-learning}

联邦学习是一种在多参与方之间协同训练模型的方法，其基本思想是将训练数据保留在本地，仅上传模型参数或梯度更新，由中心协调端完成聚合并向各节点分发新模型\cite{placeholder_federated_learning_survey,placeholder_fedavg}。对于车联网场景，这种方法能够减少原始数据跨节点流动，有助于满足“数据不出车端”的基本要求。结合当前项目实现，联邦训练流程大致可概括为：首先完成数据预处理并生成全局训练集与验证集；其次将训练集按客户端数量切分到本地文件；然后由协调端在每一轮中根据节点评分选择部分可用客户端参与训练；各客户端在本地执行若干轮梯度更新后，将更新后的参数上传；最后由服务器端执行聚合，并在验证集上评估全局模型效果。

FedAvg 是当前项目的基础聚合思路，其核心在于对各客户端本地训练后的模型参数进行按样本数加权平均。若第 $t$ 轮选中的客户端集合记为 $S_t$，第 $i$ 个客户端样本量为 $n_i$，本地训练得到的模型参数为 $w_i^{(t+1)}$，则全局模型更新可写为
\begin{equation}
    \label{eq:fedavg}
    w^{(t+1)} = \sum_{i \in S_t} \frac{n_i}{\sum_{j \in S_t} n_j} w_i^{(t+1)}
\end{equation}
在 \texttt{federated\_learning.py} 中，\texttt{aggregate\_weights()} 正是按照上述思想对浮点参数执行加权求和，因此本文后续关于聚合机制的描述均以这一实现为依据。

考虑到车联网数据往往并非独立同分布，当前实现还兼容 FedProx 形式的近端正则化策略。当 \texttt{fedprox\_mu} 大于零时，本地训练损失函数会在普通分类损失之外，增加当前模型参数与全局参考参数之间的距离惩罚，以减弱局部训练过度偏离全局目标的问题。其思想可以表示为
\begin{equation}
    \label{eq:fedprox}
    \mathcal{L}_{\text{local}} = \mathcal{L}_{\text{task}} + \frac{\mu}{2} \sum_k \left\| w_k - w_k^{(t)} \right\|_2^2
\end{equation}
代码中的 \texttt{\_compute\_fedprox\_penalty()} 即实现了这一项。虽然当前实验记录主要展示的是改进方案与标准 FedAvg 基线，但 FedProx 的接口已经在系统中预留，为后续更丰富的对照实验提供了支持。

\section{轻量化检测模型与训练优化机制}
\label{sec:related-model-optimization}

当前项目采用 `LightweightCNNLSTM` 作为入侵检测模型。该模型首先使用两层一维卷积与批归一化模块提取局部特征，再通过最大池化降低特征长度，随后把卷积输出送入 LSTM 以建模序列关联，最后利用两层全连接分类器输出预测结果。与参数规模更大的深层模型相比，该结构兼顾了特征提取能力与部署轻量化需求，适合在资源受限环境下作为车联网入侵检测器。根据现有报告文件，模型总参数量为 8402 个，模型大小约为 33608 字节，能够为后续通信统计和压缩分析提供量化基础。

为了适应车联网环境中节点状态差异明显的特点，系统在每一轮训练前会对可用客户端执行三维度动态节点选择。\texttt{score\_available\_clients()} 使用节点算力、电量水平和信道质量计算综合评分，并在此基础上引入轻量公平性惩罚，避免同一节点连续多轮被优先选中。若以 $c_i$、$b_i$、$h_i$ 分别表示第 $i$ 个客户端的算力、电量与信道质量，$w_c$、$w_b$、$w_h$ 表示对应权重，$t_i$ 表示该节点历史入选次数，则其最终评分可概括为
\begin{equation}
    \label{eq:client-selection}
    s_i = \max \left(0, \frac{w_c c_i + w_b b_i + w_h h_i}{w_c + w_b + w_h} - \min(0.05 t_i, 0.2) \right)
\end{equation}
系统随后按评分从高到低排序，并根据 \texttt{client\_fraction} 确定当轮参与节点数。该机制体现了车联网训练中“资源感知 + 公平性约束”的基本思想。

除节点筛选外，参数压缩也是降低联邦训练通信压力的重要手段。当前实现的 \texttt{\_compress\_client\_update()} 首先计算客户端模型与全局模型之间的参数增量，然后对增量向量选取绝对值最大的 Top-K 元素，仅保留最具代表性的更新分量；在此基础上，系统还支持 8/16/32 位的定点量化，其中 32 位相当于不量化，8 位和 16 位则通过缩放、离散化与反量化重建近似更新值。该策略能够在保持主要更新方向的同时显著减少实际传输字节数。对于论文分析而言，Top-K 稀疏化体现了“只发送重要参数”的思想，而定点量化体现了“用更少比特表达更新值”的思想，两者叠加构成了当前项目中的通信优化核心。

需要指出的是，当前代码中所谓“隐私代理分数”并非严格意义上的差分隐私或密码学安全证明，而是一种结合数据本地保留程度与模型更新压缩收益的代理指标。系统默认原始数据留在本地，因此基础分与数据本地性相关；同时，若压缩带来更少的上传信息量，则代理分数相应提高。该定义更适合用于工程比较，而不宜直接解释为形式化隐私保障。

\section{评价指标}
\label{sec:related-metrics}

为了从检测效果与系统开销两个层面对方案进行评价，本文采用分类性能指标与工程性能指标相结合的方式。对二分类任务而言，准确率（Accuracy）表示总体预测正确的样本比例；精确率（Precision）衡量被判为攻击的样本中真实攻击样本所占比例；召回率（Recall）反映真实攻击样本被成功识别的比例；F1 值则是精确率与召回率的调和平均。其定义分别可写为
\begin{equation}
    \label{eq:metric-accuracy}
    \mathrm{Accuracy} = \frac{TP + TN}{TP + TN + FP + FN},
\end{equation}
\begin{equation}
    \label{eq:metric-precision-recall-f1}
    \mathrm{Precision} = \frac{TP}{TP + FP}, \quad
    \mathrm{Recall} = \frac{TP}{TP + FN}, \quad
    \mathrm{F1} = \frac{2 \times \mathrm{Precision} \times \mathrm{Recall}}{\mathrm{Precision} + \mathrm{Recall}}。
\end{equation}
其中，$TP$、$TN$、$FP$、$FN$ 分别表示真正例、真反例、假正例和假反例。\texttt{metrics.py} 中的 \texttt{calculate\_classification\_metrics()} 即按这一思路计算 Accuracy、Precision、Recall 和 F1 值；当标签为二分类时，系统还会进一步根据混淆矩阵计算假正率（False Positive Rate）。

除检测性能外，联邦学习系统还必须关注通信与计算开销。当前项目通过训练历史与联邦报告记录以下指标：总训练时间、平均轮次耗时、平均本地训练耗时、原始通信量、压缩后通信量以及总体通信压缩率。总训练时间与平均轮次耗时可反映整体吞吐能力，原始通信量与压缩后通信量则直接体现参数压缩策略的收益。对于论文实验部分，这些指标将与 Accuracy、Recall 和 F1 值共同构成改进方案和基线方案的对照分析依据。

另外，系统在每轮训练后还会生成隐私代理分数，用于反映“数据保留在本地”与“上传更新量下降”这两类因素的综合效果。虽然该指标不具备严格理论证明，但在当前工程原型中能够辅助说明改进方案在减少上传信息暴露面的同时，是否仍能保持可接受的检测性能。因此，后续实验分析将采用“检测性能、通信收益、训练收益、隐私代理”四个维度进行综合讨论。

\section{本章小结}
\label{sec:related-summary}

本章围绕本文所需的关键理论基础进行了说明。首先，分析了车联网环境中的安全挑战与入侵检测任务特点，指出该问题不仅涉及分类模型设计，还受到分布式部署、资源约束和隐私要求的共同影响；其次，介绍了联邦学习的基本流程，以及 FedAvg 和 FedProx 在当前系统中的实现意义；随后，结合代码中真实存在的 `LightweightCNNLSTM`、动态节点选择、Top-K 稀疏化和定点量化机制，阐述了系统实现所依托的主要技术路线；最后，对分类指标、通信指标和隐私代理指标进行了统一说明。以上内容为第\ref{cha:system-design}章的系统分析与总体设计奠定了理论与术语基础。
